\section{Dif\/ference Galois theory: reminders and complements}
\label{sec2}

\subsection{Generalities on dif\/ference Galois theory}
%\label{sec21}

For details on what follows, we refer to~\cite[Chapter~1]{VdPS97}.

A dif\/ference ring $(R,\phi)$ is a~ring~$R$ together with a~ring automorphism $\phi: R \rightarrow R$.
An ideal of~$R$ stabilized by~$\phi$ is called a~dif\/ference ideal of $(R,\phi)$.
If~$R$ is a~f\/ield then $(R,\phi)$ is called a~dif\/ference f\/ield.

The ring of constants $R^{\phi}$ of the dif\/ference ring $(R,\phi)$ is def\/ined~by
\begin{gather*}
R^{\phi}:=\{f \in R \, | \, \phi(f)=f\}.
\end{gather*}

A dif\/ference ring morphism (resp.\
dif\/ference ring isomorphism) from the dif\/ference ring $(R,\phi)$ to the dif\/ference ring
$(\widetilde{R},\widetilde{\phi})$ is a~ring morphism (resp.\
ring isomorphism) $\varphi: R \rightarrow \widetilde{R}$ such that $\varphi \circ \phi = \widetilde{\phi} \circ
\varphi$.

A dif\/ference ring $(\widetilde{R},\widetilde{\phi})$ is a~dif\/ference ring extension of a~dif\/ference ring $(R,\phi)$ if
$\widetilde{R}$ is a~ring extension of~$R$ and $\widetilde{\phi}_{\vert R}=\phi$; in this case, we will often denote
$\widetilde{\phi}$ by~$\phi$.
Two dif\/ference ring extensions $(\widetilde{R}_{1},\widetilde{\phi}_{1})$ and $(\widetilde{R}_{2},\widetilde{\phi}_{2})$
of a~dif\/ference ring $(R,\phi)$ are isomorphic over $(R,\phi)$ if there exists a~dif\/ference ring isomorphism~$\varphi$
from $(\widetilde{R}_{1},\widetilde{\phi}_{1})$ to $(\widetilde{R}_{2},\widetilde{\phi}_{2})$ such that $\varphi_{\vert
R}=\operatorname{Id}_{R}$.

{\it We now let $(K,\phi)$ be a~difference field.
We assume that its field of constants $C:=K^{\phi}$ is algebraically closed and that the characteristic of~$K$ is~$0$.}

Consider a~linear dif\/ference system
\begin{gather}
\label{the generic system}
\phi Y=AY \qquad \hbox{with} \quad A \in \operatorname{GL}_{n}(K).
\end{gather}

According to~\cite[Section~1.1]{VdPS97}, there exists a~dif\/ference ring extension $(R,\phi)$ of $(K,\phi)$ such that
\begin{itemize}\itemsep=0pt
\item[1)] there exists $U \in \operatorname{GL}_{n}(R)$ such that $\phi (U) = AU$ (such a~$U$ is called a~fundamental
matrix of solutions of~\eqref{the generic system});
\item[2)] $R$ is generated, as a~$K$-algebra, by the entries of~$U$ and $\det(U)^{-1}$;
\item[3)] the only dif\/ference ideals of $(R,\phi)$ are $\{0\}$ and~$R$.
\end{itemize}
Such a~dif\/ference ring $(R,\phi)$ is called a~Picard--Vessiot ring for~\eqref{the generic system} over~$(K,\phi)$.
It is unique up to isomorphism of dif\/ference rings over $(K,\phi)$.
It is worth mentioning that $R^{\phi}=C$; see~\cite[Lemma~1.8]{VdPS97}.
\begin{Remark}
Picard--Vessiot rings are not domains in general: they are f\/inite direct sums of domains cyclically permuted by~$\phi$;
see~\cite[Corollary~1.16]{VdPS97}.
\end{Remark}

The corresponding dif\/ference Galois group~$G$ over $(K,\phi)$ of~\eqref{the generic system} is the group of~$K$-linear
ring automorphisms of~$R$ commuting with~$\phi$:
\begin{gather*}
G:=\{\sigma \in \operatorname{Aut}(R/K) \, | \, \phi\circ\sigma=\sigma\circ \phi \}.
\end{gather*}
The choice of the base f\/ield is by no way innocent.
The bigger the base f\/ield is, the smaller the Galois group is.

A straightforward computation shows that, for any~$\sigma {\in} G$, there exists a~unique ${C(\sigma) {\in}
\operatorname{GL}_{n}(C)}$ such that $\sigma(U)=UC(\sigma)$.
According to~\cite[Theorem~1.13]{VdPS97}, one can identify~$G$ with an {\it algebraic} subgroup of
$\operatorname{GL}_{n}(C)$ via the faithful representation
\begin{gather*}
\sigma \in G \mapsto C(\sigma) \in \operatorname{GL}_{n}(C).
\end{gather*}
If we choose another fundamental matrix of solutions~$U$, we f\/ind a~conjugate representation.

\begin{Remark}
Given an~$n$th order dif\/ference equation
\begin{gather}
\label{the generic equa}
a_{n} \phi^{n} (y) + \cdots + a_{1} \phi(y)+a_{0} y=0,
\end{gather}
with $a_{0},\dots,a_{n} \in K$ and $a_{0}a_{n} \neq 0$, we can consider the equivalent linear dif\/ference system
\begin{gather}
\label{matrix system equa}
\phi Y=AY, \qquad \text{with} \quad  A=
\begin{pmatrix}
0&1&0&\dots&0
\\
0&0&1&\ddots&\vdots
\\
\vdots&\vdots&\ddots&\ddots&0
\\
0&0&\dots&0&1
\\
-\frac{a_{0}}{a_{n}}& -\frac{a_{1}}{a_{n}}&\dots & \dots & -\frac{a_{n-1}}{a_{n}}
\end{pmatrix}
\in \operatorname{GL}_{n}(K).
\end{gather}
By ``Galois group of the dif\/ference equation~\eqref{the generic equa}'' we mean ``Galois group of the dif\/ference
system~\eqref{matrix system equa}''.
\end{Remark}

We shall now introduce a~property relative to the dif\/ference base f\/ield, which appears in~\cite[Lemma~1.19]{VdPS97}.

\begin{Definition}
\label{prop P}
We say that the dif\/ference f\/ield $(K,\phi)$ satisf\/ies property $(\mathcal{P})$ if the following properties hold:
\begin{itemize}\itemsep=0pt
\item the f\/ield~$K$ is a~$\mathcal{C}^{1}$-f\/ield\footnote{Recall that~$K$ is a~$\mathcal{C}^{1}$-f\/ield if every
non-constant homogeneous polynomial~$P$ over~$K$ has a~non-trivial zero provided that the number of its variables is
more than its degree.};
\item if~$L$ is a~f\/inite f\/ield extension of~$K$ such that~$\phi$ extends to a~f\/ield endomorphism of~$L$ then $L=K$.
\end{itemize}
\end{Definition}

The following result is due to van der Put and Singer.
We recall that two dif\/ference systems $\phi Y=AY$ and $\phi Y=BY$ with $A,B \in \operatorname{GL}_{n}(K)$ are isomorphic
over~$K$ if and only if there exists $T \in \operatorname{GL}_{n}(K)$ such that $\phi(T) A=BT$.

\begin{Theorem}
\label{si prop P alors}
Assume that $(K,\phi)$ satisfies property $(\mathcal{P})$.
Let $K^{\phi} = C$.
Let $G \subset \operatorname{GL}_{n}(C)$ be the difference Galois group over $(K,\phi)$ of
\begin{gather}
\label{the generic system theo}
\phi(Y) = AY, \qquad \text{with} \quad  A \in \operatorname{GL}_{n}(K).
\end{gather}
Then, the following properties hold:
\begin{itemize}
\itemsep=0pt
\item $G/G^{\circ}$ is cyclic, where $G^{\circ}$ is the identity component of~$G$;
\item there exists $B \in G(K)$ such that~\eqref{the generic system theo} is isomorphic to $\phi Y=BY$ over~$K$.
\end{itemize}
Let $\widetilde{G}$ be an algebraic subgroup of $\operatorname{GL}_{n}(C)$ such that~$A\in \widetilde{G}(K)$.
The following properties hold:
\begin{itemize}
\itemsep=0pt
\item~$G$ is conjugate to a~subgroup of $\widetilde{G}$;
\item any minimal element in the set of algebraic subgroups $\widetilde{H}$ of $\widetilde{G}$ for which there exists
$T\in \operatorname{GL}_{n}(K)$ such that $\phi(T)AT^{-1}\in \widetilde{H}(K)$ is conjugate to~$G$;
\item~$G$ is conjugate to $\widetilde{G}$ if and only if, for any $T\in \widetilde{G}(K)$ and for any proper algebraic
sub\-group~$\widetilde{H}$ of~$\widetilde{G}$, one has that $\phi(T)AT^{-1}\notin \widetilde{H}(K)$.
\end{itemize}
\end{Theorem}

\begin{proof}
The proof of~\cite[Propositions~1.20 and~1.21]{VdPS97} in the special case where $K:=C(z)$ and~$\phi$ is the shift
$\phi(f(z)):=f(z+h)$ with $h \in C^\times$, extends {\it mutatis mutandis} to the present case.
\end{proof}

\subsection{Dif\/ference equations on elliptic curves}

Let $\Lambda \subset \mathbb C$ be a~lattice.
Without loss of generality, we can assume that
\begin{gather*}
\Lambda=\mathbb Z +\mathbb Z \tau, \qquad \text{with} \quad \Im (\tau)>0,
\end{gather*}
where $\Im(\cdot)$ denotes the imaginary part.
For any lattice $\Lambda' \subset \mathbb C$, we let $\operatorname{M}_{\Lambda'}$ be the f\/ield of $\Lambda'$-periodic
meromorphic functions.
We denote by~$K$ the f\/ield def\/ined by
\begin{gather*}
K:= \bigcup_{\Lambda' \; \text{sub-lattice of} \; \Lambda} \operatorname{M}_{\Lambda'} =\bigcup_{k \geq 1} \operatorname{M}_{k  \Lambda}.
\end{gather*}
Let $h \in \mathbb C$ such that $h \mod \Lambda$ is not a~torsion point of $\mathbb C/\Lambda$.
We endow~$K$ with the non-cyclic f\/ield automorphism~$\phi$ def\/ined~by
\begin{gather*}
\phi(f)(z):=f(z+h).
\end{gather*}
Then,~$(K,\phi)$ is a~dif\/ference f\/ield.

\begin{Proposition}
The field of constants of $(K,\phi)$ is
\begin{gather*}
K^{\phi}=\mathbb C.
\end{gather*}
\end{Proposition}

\begin{proof}
Consider $f\in K^{\phi}$.
Let $\Lambda'$ be a~sub-lattice of~$\Lambda$ such that $f \in \operatorname{M}_{\Lambda'}$.
Note that~$f$ is $\Lambda'$-periodic (because $f \in \operatorname{M}_{\Lambda'}$) and~$h$-periodic (because
$\phi(f)=f$), so~$f$ is a~$(\Lambda' + h\mathbb Z)$-periodic meromorphic function.
But $\Lambda' + h\mathbb Z$ has an accumulation point because $h \mod \Lambda$ is not a~torsion point of $\mathbb
C/\Lambda$.
Therefore,~$f$ is constant.
\end{proof}

\begin{Proposition}
The difference field $(K,\phi)$ satisfies property $(\mathcal{P})$ $($see Definition~{\rm \ref{prop P})}.
\end{Proposition}

\begin{proof}
Since $K = \bigcup_{k \geq 1} \operatorname{M}_{k \Lambda}$ is the increasing union of the f\/ields
$\operatorname{M}_{k \Lambda}$, the fact that~$K$ is a~$\mathcal{C}^{1}$-f\/ield follows from Tsen's
theorem~\cite{LangQAC} (according to which the function f\/ield of any algebraic curve over an algebraically closed f\/ield,
e.g.,
$\operatorname{M}_{k \Lambda}$, is $\mathcal{C}^{1}$).

Let~$L$ be a~f\/inite extension of~$K$ such that~$\phi$ extends to a~f\/ield endomorphism of~$L$.
We have to prove that $L=K$.
The primitive element theorem ensures that there exists $u \in L$ such that $L=K(u)$.
Let $\Lambda'$ be a~sub-lattice of~$\Lambda$ such that
\begin{itemize}
\itemsep=0pt
\item~$u$ is algebraic over $\operatorname{M}_{\Lambda'}$,
\item $\phi(u) \in \operatorname{M}_{\Lambda'}(u)$.
\end{itemize}
Then, $\operatorname{M}_{\Lambda'}(u)$ is a~f\/inite extension of $\operatorname{M}_{\Lambda'}$ and~$\phi$ induces an
automorphism of $\operatorname{M}_{\Lambda'}(u)$.

Using the equivalence of categories between between smooth projective curves and function f\/ields of dimension
$1$~\cite[Corollary 6.12]{HartshorneAG}, we see that there exists a~commutative diagram of the form
\begin{gather*}
\xymatrix{X \ar[r]^{f} \ar[d]_{\varphi} &  X \ar[d]^{\varphi}
\\
\mathbb C/\Lambda' \ar[r]_{z \mapsto z+h}  & \mathbb C/\Lambda'}
\end{gather*}
where $\varphi: X \rightarrow \mathbb C/\Lambda'$ is a~morphism of smooth projective curves, whose induced morphism of
function f\/ields ``is'' the inclusion $\operatorname{M}_{\Lambda'} \subset \operatorname{M}_{\Lambda'}(u)$, and where~$f$
is an endomorphism of~$X$, whose induced morphism on function f\/ields ``is'' $\phi: \operatorname{M}_{\Lambda'}(u)
\rightarrow \operatorname{M}_{\Lambda'}(u)$.
Considering this commutative diagram, we see that~$f$ has degree $1$ and that, if~$\varphi$ is ramif\/ied above $y \in
\mathbb C/\Lambda'$, then~$\varphi$ is also ramif\/ied above $y-h$.
So, the set of ramif\/ication values of~$\varphi$ is stable by $z \mapsto z-h$.
This set being f\/inite, it has to be empty.
So,~$\varphi$ is unramif\/ied.

Hurwitz's formula implies that~$X$ has genus~$1$, i.e., that~$X$ is an elliptic curve.
So, there exist a~lattice $\Lambda'' \subset \mathbb C$ and an isomorphism $\psi: \operatorname{M}_{\Lambda'}(u)
\rightarrow \operatorname{M}_{\Lambda''}$.
There exists $(a,b) \in \mathbb C^{\times} \times \mathbb C$ such that $a \Lambda'' \subset \Lambda'$ and such that the
restriction $\psi_{\vert \operatorname{M}_{\Lambda'}}$ is given, for all $f \in \operatorname{M}_{\Lambda'}$,~by
$\psi(f)(z)=f(az+b)$.
(Indeed, $\psi_{\vert \operatorname{M}_{\Lambda'}}: \operatorname{M}_{\Lambda'} \rightarrow
\operatorname{M}_{\Lambda''}$ is a~f\/ield morphism from the function f\/ield of the elliptic curve~$\mathbb C/\Lambda'$ to
the function f\/ield of the elliptic curve~$\mathbb C/\Lambda''$.
So, $\psi_{\vert \operatorname{M}_{\Lambda'}}$ is induced by a~morphism from the elliptic curve $\mathbb C/\Lambda''$ to
the elliptic curve $\mathbb C/\Lambda'$.
Now, our claim follows from the fact that the morphisms from $\mathbb C/\Lambda''$ to $\mathbb C/\Lambda'$ are of the
form $z \mod \Lambda'' \mapsto az+b \mod \Lambda'$ for some $(a,b) \in \mathbb C^{\times} \times \mathbb C$ such that $a
\Lambda'' \subset \Lambda'$.) The commutative diagram
\begin{gather*}
\xymatrix{\operatorname{M}_{\Lambda'} \ar@{^{(}->}[r] \ar[rd]_{v(z) \mapsto v(az+b)} &  \operatorname{M}_{\Lambda'}(u)
\ar[d]^{\psi} \ar[rd] &
\\
&  \operatorname{M}_{\Lambda''} \ar[r]_{w(z) \mapsto w(\frac{z-b}{a})}  & \operatorname{M}_{a \Lambda''}}
\end{gather*}
shows that the f\/ields $\operatorname{M}_{\Lambda'}(u)$ and $\operatorname{M}_{a \Lambda''}$ are
$\operatorname{M}_{\Lambda'}$-isomorphic.
But the extension $\operatorname{M}_{a \Lambda''}/\operatorname{M}_{\Lambda'}$ is Galois (indeed, this is equivalent to
the fact that the corresponding morphism of smooth projective curves $\mathbb C/\Lambda' \rightarrow \mathbb
C/a\Lambda''$ is Galois, and this is easily seen from the explicit description of the morphisms between these curves).
Therefore, any $\operatorname{M}_{\Lambda'}$-morphism from $\operatorname{M}_{a \Lambda''}$ to $K (u)$ must leave
$\operatorname{M}_{a \Lambda''}$ globally invariant.
But, $\operatorname{M}_{a \Lambda''}$ and $\operatorname{M}_{\Lambda'}(u)$ are $\operatorname{M}_{\Lambda'}$-isomorphic
subf\/ields of~$K(u)$.
So $\operatorname{M}_{\Lambda'} (u) \subset \operatorname{M}_{a \Lambda''}$, and therefore $u \in \operatorname{M}_{a
\Lambda''} \subset K$ and $L=K(u) \subset K$.
\end{proof}

\begin{Corollary}
\label{coro1}
The conclusions of Theorem~{\rm \ref{si prop P alors}} are valid for $(K,\phi)$.
\end{Corollary}
